% ===== PRÉAMBULE CANONIQUE — à coller VERBATIM en tête de CHAQUE .tex =====
\documentclass[11pt,a4paper]{article}
\usepackage[utf8]{inputenc}\usepackage[T1]{fontenc}\usepackage{lmodern}
\usepackage[french]{babel}
\usepackage{amsmath,amssymb,mathtools}\usepackage{icomma}
\usepackage[version=4]{mhchem}          % \ce{...} : équations & formules
\usepackage{siunitx}\sisetup{locale=FR} % \qty{}{}, \si{}, \ang{}
\usepackage{chemfig}                    % \chemfig{...} : structures développées
\usepackage{xcolor}\usepackage{colortbl}
\usepackage{tikz}\usepackage{pgfplots}\pgfplotsset{compat=1.18}
\usepackage[most]{tcolorbox}\usepackage{enumitem}\usepackage{array,booktabs}
\usepackage[a4paper,margin=2.2cm]{geometry}
\usetikzlibrary{arrows.meta,calc,angles,quotes,positioning,patterns,backgrounds,shapes.geometric}
\definecolor{primary}{RGB}{222,49,99}\definecolor{primarydark}{RGB}{150,22,55}
\definecolor{defcol}{RGB}{0,114,178}\definecolor{methcol}{RGB}{52,92,148}
\definecolor{excol}{RGB}{0,135,90}\definecolor{attncol}{RGB}{200,40,55}
\tcbset{boxbase/.style={enhanced,breakable,boxrule=0.9pt,arc=2.5pt,
  left=3mm,right=3mm,top=2.3mm,bottom=2.3mm,fonttitle=\bfseries,coltitle=white}}
% --- boîtes TUTORIEL (noms EXACTS, ne pas renommer) ---
\newtcolorbox{cle}[1][]{boxbase,colback=defcol!6,colframe=defcol,
  title={\textsc{À retenir}\ifx\relax#1\relax\relax\else\textmd{ : \itshape #1}\fi}}
\newtcolorbox{methode}[1][]{boxbase,colback=methcol!7,colframe=methcol,sharp corners,
  title={\textsc{Méthode}\ifx\relax#1\relax\relax\else\textmd{ --- \itshape #1}\fi}}
\newtcolorbox{exemplebox}[1][]{boxbase,colback=excol!5,colframe=excol,
  title={\textsc{Exemple}\ifx\relax#1\relax\relax\else\textmd{ : \itshape #1}\fi}}
\newtcolorbox{piege}[1][]{boxbase,colback=attncol!6,colframe=attncol,
  title={\textsc{Piège}\ifx\relax#1\relax\relax\else\textmd{ : \itshape #1}\fi}}
\newtcolorbox{remarque}[1][]{boxbase,colback=black!4,colframe=black!45,
  title={\textsc{Remarque}\ifx\relax#1\relax\relax\else\textmd{ : \itshape #1}\fi}}
% --- boîtes EXERCICE / CORRIGÉ (noms EXACTS, auto-comptées) ---
\newtcolorbox[auto counter]{exo}[1][]{boxbase,colback=primary!4,colframe=primary,
  title={\textsc{Exercice}~\thetcbcounter\ifx\relax#1\relax\relax\else\textmd{ --- \itshape #1}\fi}}
\newtcolorbox[auto counter]{corrige}[1][]{boxbase,colback=excol!4,colframe=excol,
  title={\textsc{Corrigé}~\thetcbcounter\ifx\relax#1\relax\relax\else\textmd{ --- \itshape #1}\fi}}
\newcommand{\R}{\mathbb{R}}\newcommand{\N}{\mathbb{N}}\newcommand{\Z}{\mathbb{Z}}\newcommand{\ds}{\displaystyle}
% (un \usepackage{fancyhdr}+en-tête est possible ; il est ignoré côté web)
% =========================================================================
\begin{document}
\begin{center}\color{primarydark}\large\bfseries Thème 1 — Nombre d'oxydation \quad$\bullet$\quad Niveau Difficile\end{center}

\begin{exo}[vrai ou faux : les définitions]
Pour chaque proposition, indique si elle est \textbf{vraie} ou \textbf{fausse} (et corrige les
fausses). Combien de propositions sont vraies ?
\begin{enumerate}[label=\alph*)]
  \item Le nombre d'oxydation d'un ion monoatomique est égal à sa charge.
  \item Lorsqu'un élément s'oxyde, son nombre d'oxydation diminue.
  \item L'oxydant est l'espèce qui capte des électrons et qui est réduite au cours de la réaction.
  \item Le nombre d'oxydation de l'hydrogène vaut toujours $+\mathrm{I}$.
  \item Dans un corps simple, le nombre d'oxydation d'un atome est nul.
\end{enumerate}
\end{exo}

\begin{exo}[tous les éléments d'un composé, et un piège]
On s'intéresse au permanganate de potassium \ce{KMnO4}.
\begin{enumerate}[label=\alph*)]
  \item Détermine le nombre d'oxydation de \textbf{chaque} élément : \ce{K}, \ce{Mn} et \ce{O}.
  \item Un étudiant écrit $\mathrm{n.o.}(\ce{Mn})=+\mathrm{VIII}$. Explique son erreur et donne la
        bonne valeur.
\end{enumerate}
\end{exo}

\begin{exo}[reconnaître une réaction d'oxydo-réduction]
Parmi les réactions suivantes, lesquelles sont des réactions d'\textbf{oxydo-réduction} ? Justifie en
repérant les atomes dont le nombre d'oxydation \textbf{varie}.
\begin{enumerate}[label=\alph*)]
  \item \ce{NaCl + AgNO3 -> AgCl v + NaNO3}
  \item \ce{Cu + 4 HNO3 -> Cu(NO3)2 + 2 NO2 ^ + 2 H2O}
  \item \ce{2 Na + 2 H2O -> 2 NaOH + H2 ^}
  \item \ce{CO2 + H2O -> H2CO3}
\end{enumerate}
\end{exo}

\begin{exo}[un élément oxydé \emph{et} réduit à la fois]
On considère la réaction (non équilibrée)
\[
\ce{MnSO4 + KMnO4 + H2O -> MnO2 + H2SO4 + K2SO4}.
\]
\begin{enumerate}[label=\alph*)]
  \item Détermine le nombre d'oxydation du manganèse dans \ce{MnSO4}, \ce{KMnO4} et \ce{MnO2}.
  \item Quel manganèse est \textbf{oxydé} ? Lequel est \textbf{réduit} ?
  \item Le permanganate \ce{KMnO4} joue-t-il le rôle d'oxydant ou de réducteur ?
\end{enumerate}
\end{exo}

\begin{exo}[relier n.o., formule et sens de la réaction]
Les sels de fer usuels sont le sulfate de fer(II) \ce{FeSO4} et le sulfate de fer(III)
\ce{Fe2(SO4)3} (l'ion sulfate est \ce{SO4^{2-}}).
\begin{enumerate}[label=\alph*)]
  \item Quel est le nombre d'oxydation du fer dans \ce{FeSO4} ?
  \item Quel est le nombre d'oxydation du fer dans \ce{Fe2(SO4)3} ?
  \item Lorsqu'on transforme \ce{FeSO4} en \ce{Fe2(SO4)3}, le fer est-il oxydé ou réduit ? Combien
        d'électrons chaque atome de fer échange-t-il ?
\end{enumerate}
\end{exo}

\end{document}
